\documentclass[../main.tex]{subfiles}
\begin{document}

\chapter{ILP}
\lessoninfo{
  video={Lesson 6, videos 1--20},
  textbook={H\&P \cite{hennessy2017} §3.1, §3.4; \cite{keller1975,wall1991}},
  keywords={RAW dependence, WAW dependence, false dependence, architectural register, physical register, register renaming, register alias table, instruction-level parallelism, ideal processor, instructions per cycle, processor width, in-order execution, out-of-order execution},
}

Branch prediction (Lesson 4) and predication (Lesson 5) remove most of the
cycles that control dependences cost a pipeline, which brings its CPI close
to the ideal value of~1.  Going below~1 requires executing several
instructions in the same cycle, and what prevents this are the data
dependences between instructions.  This lesson shows that only
read-after-write dependences are fundamental, that the others are removed by
register renaming, and that the remaining dependences define the
instruction-level parallelism of a program: the bound against which the
instructions per cycle of real processors are measured.

\section{All instructions in the same cycle}
\label{sec:l06-same-cycle}

A pipeline that never stalls or flushes completes one instruction per cycle,
so its CPI is~1 (\cref{eq:pipeline-cpi}).\source{Lesson 6, videos 1--3; H\&P
§3.1.}  A lower CPI means that more than one instruction completes per
cycle.  In the extreme, all $N$ instructions of a program would be fetched in
the same cycle, decoded together in the next cycle, executed together in the
cycle after that, and so on (\cref{fig:l06-same-cycle}).  A five-stage
processor would complete the whole program in 5 cycles, a CPI of $5/N$,
which tends to~0 for long programs.

A real program cannot be executed this way, because its instructions use
the results of other instructions.  Consider the following five
instructions, the running example of the first half of this lesson:%
\begin{marginfigure}
  \resizebox{\linewidth}{!}{% Thought experiment: all N instructions pass through the five stages in
% the same five cycles.
\begin{tikzpicture}[x=0.9cm,y=0.55cm, font=\footnotesize\sffamily,
  st/.style={draw, minimum width=0.82cm, minimum height=0.46cm, inner sep=0pt, fill=ln-box}]
  \foreach \c in {1,...,5} { \node at (\c,0.8) {\c}; }
  \node[anchor=east] at (0.5,0.8) {\iflangja{サイクル}{cycle}};
  \foreach \r/\name in {0/I1,1/I2,2/I3,4/I$N$} {
    \node[anchor=east] at (0.5,-\r) {\name};
    \foreach \s/\lab in {1/IF,2/ID,3/EX,4/MEM,5/WB} {
      \node[st] at (\s,-\r) {\lab};
    }
  }
  \foreach \s in {1,...,5} { \node at (\s,-3) {$\vdots$}; }
\end{tikzpicture}
}
  \caption{The ideal: all $N$ instructions pass through the stages in the
    same five cycles.}
  \label{fig:l06-same-cycle}
\end{marginfigure}
\begin{lstlisting}[language={[x86masm]Assembler}, numbers=none]
I1:  ADD R3, R1, R2     ; R3 <- R1 + R2
I2:  SUB R5, R3, R4     ; R5 <- R3 - R4
I3:  AND R6, R1, R4     ; R6 <- R1 AND R4
I4:  OR  R4, R2, R7     ; R4 <- R2 OR R7
I5:  XOR R5, R1, R7     ; R5 <- R1 XOR R7
\end{lstlisting}
When all five instructions are in the execute stage in the same cycle, I2
computes $\text{R3} - \text{R4}$ while I1 is still computing R3, so I2 uses
the old value of R3.  Forwarding (\cref{fig:forwarding}) cannot correct
this: it delivers a result computed in one cycle to an instruction that
executes in a later cycle, whereas here the result and its use fall into the
same cycle.  How close a processor can come to the ideal is therefore
determined by the data dependences of the program.

\section{Dependences in same-cycle execution}
\label{sec:l06-dependences}

The read-after-write dependence of I2 on I1 must be respected: I2 can execute
no earlier than one cycle after I1, when the result of I1 can be forwarded to
it.\source{Lesson 6, videos 4--6; H\&P §3.1.}  I3, I4 and I5 read no register
that an earlier instruction writes, so they can still execute in the same
cycle as I1.  \Cref{fig:l06-dependences} shows the resulting timing: I2 waits
one cycle, and I3 to I5 execute before I2 although they follow it in the
program.  RAW dependences therefore bound how far the CPI can fall even with
unlimited width: if every instruction read the result of its predecessor,
one instruction would complete per cycle and the CPI would remain~1.

\begin{figure}[htbp]
  \centering
  % Same-cycle execution of the five-instruction program: the RAW dependence
% delays I2, which then overwrites the later write of I5 to R5 (WAW).
\begin{tikzpicture}[x=1.2cm,y=0.9cm, font=\footnotesize\sffamily, >={Stealth[length=1.8mm]},
  st/.style={draw, minimum width=0.95cm, minimum height=0.5cm, inner sep=0pt, fill=ln-box},
  bad/.style={st, fill=ln-caution!15, draw=ln-caution},
  stall/.style={draw=ln-rule, dashed, minimum width=0.95cm, minimum height=0.5cm,
    inner sep=0pt, font=\scriptsize\sffamily, text=ln-muted}]
  \foreach \c in {1,...,6} { \node at (\c,0.7) {\c}; }
  \node[anchor=east] at (0.45,0.7) {\iflangja{サイクル}{cycle}};
  \foreach \r/\txt in {1/{ADD R3,R1,R2},2/{SUB R5,R3,R4},3/{AND R6,R1,R4},
                       4/{OR~~R4,R2,R7},5/{XOR R5,R1,R7}} {
    \node[anchor=east] at (0.45,1-\r) {I\r\enspace\texttt{\txt}};
  }
  % I1, I3, I4: undelayed
  \foreach \s/\lab in {1/IF,2/ID,4/MEM,5/WB} { \node[st] at (\s,0) {\lab}; }
  \node[st] (e1) at (3,0) {EX};
  \foreach \r in {3,4} {
    \foreach \s/\lab in {1/IF,2/ID,3/EX,4/MEM,5/WB} { \node[st] at (\s,1-\r) {\lab}; }
  }
  % I2: waits one cycle for R3
  \node[st] at (1,-1) {IF};
  \node[st] at (2,-1) {ID};
  \node[stall] at (3,-1) {\iflangja{待ち}{wait}};
  \node[st] (e2) at (4,-1) {EX};
  \node[st] at (5,-1) {MEM};
  \node[bad] (w2) at (6,-1) {WB};
  % I5: independent, but writes R5 before I2 does
  \foreach \s/\lab in {1/IF,2/ID,3/EX,4/MEM} { \node[st] at (\s,-4) {\lab}; }
  \node[bad] (w5) at (5,-4) {WB};
  % forwarding of R3 from I1 to I2
  \draw[->, thick, ln-accent] (e1.south) -- (3,-0.5)
    -- node[fill=white, inner sep=0.6pt, font=\scriptsize\sffamily] {R3} (4,-0.5)
    -- (e2.north);
  % WAW violation
  \draw[->, dashed, thick, ln-caution] (w2.south) -- (w5.north east);
  \node[anchor=west, font=\scriptsize\sffamily, text=ln-caution, align=left] at (6.0,-2.5)
    {\iflangja{R5 を\\上書き}{overwrites\\R5}};
\end{tikzpicture}

  \caption{Same-cycle execution of the running example.  The RAW dependence
    on R3 delays I2 by one cycle; as a consequence I2 writes R5 after I5,
    which violates their WAW dependence.}
  \label{fig:l06-dependences}
\end{figure}

Executing instructions ahead of an earlier one breaks the write-after-write
dependence between I2 and I5.  Both write R5, and in program order I5 writes
last, so after these instructions R5 must hold the result of I5; every later
instruction that reads R5 expects that value.  In \cref{fig:l06-dependences},
however, I5 writes R5 in cycle~5 and the delayed I2 overwrites it in
cycle~6.  To keep the correct value, I5 would also have to be delayed until
after the write of I2, although its computation needs nothing from I2.

A write-after-read dependence is at risk in the same way.  I4 writes R4,
which I2 reads, and program order requires I2 to use the old value of R4.
Once I2 is delayed, the new value produced by I4 already exists when I2
executes, and I2 must be kept from using it.  A single RAW delay thus
spreads: every later instruction that writes a register the delayed
instruction reads or writes must be delayed as well, so that every
instruction still sees the register values of its position in the program.

\section{Removing false dependences}
\label{sec:l06-false}

The three kinds of data dependence of Lesson 3 differ in
nature.\source{Lesson 6, videos 7--8; H\&P §3.1.}  In a RAW dependence the
later instruction uses the value computed by the earlier one: the data
really flows from one instruction to the other, and no hardware can let the
reader execute before the writer.\index{read-after-write dependence}  WAW and
WAR dependences\index{write-after-write dependence}%
\index{write-after-read dependence} are false
dependences\index{false dependences}: no value
passes between the two instructions.  They exist only because the program
stores different values in the same register, a consequence of the limited
number of registers and of the fact that an instruction executed
repeatedly, for example in every iteration of a loop, always writes the
same register.  In the example, the
results of I2 and I5 are unrelated values that happen to be named R5.

False dependences therefore disappear if different values are kept in
different places.  One way to achieve this is to keep several versions of a
register.  When I2 and I5 both write R5, the processor stores both results
instead of letting one overwrite the other, and each instruction that reads
R5 uses the version selected by program order: an instruction located
between I2 and I5 uses the result of I2, an instruction after I5 the result
of I5.  The order in which the two writes actually happen no longer matters,
so I5 can execute as early as if I2 did not exist.  Likewise, I2 reads the
old version of R4 no matter when I4 produces the new one.

The difficulty is the bookkeeping.  For every read, the processor must find
the right version, the one produced by the closest preceding writer in
program order, and it must recognize when a version is no longer needed.
Register renaming performs this bookkeeping systematically.

\section{Register renaming}
\label{sec:l06-renaming}

Renaming separates the registers that a program names from the storage in
which the processor keeps their values.\source{Lesson 6, videos 9--11; H\&P
§3.1, §3.4; \cite{keller1975}.}

\begin{definition}[Architectural and physical registers]
  \label{def:l06-registers}
  A register defined by the instruction set, and therefore named in
  programs (for example R0 to R31), is an \term{architectural
  register}[アーキテクチャレジスタ].  A storage location in which the
  processor actually keeps a register value is a \term{physical
  register}[物理レジスタ].  A processor that renames has many more
  physical registers than architectural ones.
\end{definition}

\begin{definition}[Register renaming]
  \label{def:l06-renaming}
  Rewriting each instruction, as it is decoded, so that it reads and writes
  physical instead of architectural registers, with a newly allocated
  physical register for every result, is called \term{register
  renaming}[レジスタリネーミング].
\end{definition}

\begin{definition}[Register alias table]
  \label{def:l06-rat}
  The \term{register alias table}[レジスタエイリアステーブル]%
  \index{RAT|see{register alias table}} (RAT) has one entry per
  architectural register.  The entry names the physical register that holds
  the most recent value of that architectural register, in program order,
  among the instructions renamed so far.
\end{definition}

Each instruction is renamed in three steps (\cref{fig:l06-rat}):%
\margin{The lecture expands RAT as \emph{register allocation table}; both
names denote the same table.}
\begin{enumerate}
  \item For each source register, read its RAT entry; the renamed
    instruction reads that physical register.
  \item Allocate a free physical register, one that holds no value still
    needed, for the result; the renamed instruction writes it.
  \item Overwrite the RAT entry of the destination register with the new
    physical register, so that later instructions reading this architectural
    register use the new result.
\end{enumerate}
The sources must be looked up before the destination entry is overwritten;
otherwise an instruction such as \texttt{ADD R1, R1, R2} would read its own
result instead of the previous value of R1.\margin{Lessons 7 and~8 show how
renaming is combined with out-of-order execution in real designs.}

\begin{figure}[htbp]
  \centering
  % Renaming one instruction: sources are looked up in the RAT, the result
% gets a free physical register, and the RAT entry of the destination is
% updated to it.
\begin{tikzpicture}[x=1cm,y=1cm, font=\footnotesize\sffamily, >={Stealth[length=1.8mm]},
  fld/.style={draw, fill=ln-box, minimum width=1.1cm, minimum height=0.5cm, inner sep=0pt},
  arch/.style={draw, fill=ln-box, minimum width=0.8cm, minimum height=0.5cm, inner sep=0pt},
  phys/.style={draw, fill=white, minimum width=0.8cm, minimum height=0.5cm, inner sep=0pt}]
  % decoded instruction
  \node[fld] (s1) at (0,0)    {R1};
  \node[fld] (s2) at (0,-0.5) {R2};
  \node[fld] (d)  at (0,-1.0) {R3};
  \foreach \n/\lab in {s1/src$_1$,s2/src$_2$,d/dst} {
    \node[anchor=east, font=\scriptsize, text=ln-muted] at (\n.west) {\lab};
  }
  \node[anchor=south] at (0,0.3) {\texttt{ADD R3, R1, R2}};
  % RAT
  \foreach \i in {1,...,7} {
    \node[arch] (a\i) at (3.0,{-(\i-1)*0.5}) {R\i};
  }
  \foreach \i in {1,2,4,5,6,7} {
    \node[phys] (p\i) at (3.8,{-(\i-1)*0.5}) {P\i};
  }
  \node[phys, fill=ln-accent!15] (p3) at (3.8,-1.0) {P8};
  \node[anchor=north, inner sep=1pt] at (3.4,-3.27) {$\vdots$};
  \node[anchor=north, font=\footnotesize\bfseries\sffamily] at (3.4,-3.75) {RAT};
  % renamed instruction
  \node[fld] (r1) at (6.8,0)    {P1};
  \node[fld] (r2) at (6.8,-0.5) {P2};
  \node[fld] (rd) at (6.8,-1.0) {P8};
  \node[anchor=south] at (6.8,0.3) {\texttt{ADD P8, P1, P2}};
  % lookups
  \draw[->] (s1.east) -- (a1.west);
  \draw[->] (s2.east) -- (a2.west);
  \draw[->] (p1.east) -- (r1.west);
  \draw[->] (p2.east) -- (r2.west);
  % destination
  \draw[->, dashed] (d.east) -- (a3.west)
    node[midway, below, font=\scriptsize] {\iflangja{エントリ選択}{select entry}};
  \node[draw, fill=ln-box, minimum height=0.5cm, inner sep=3pt] (free) at (6.8,-2.4)
    {\iflangja{空き}{free}: P8, P9, P10, \ldots};
  \draw[->, thick, ln-accent] (free.north) -- (rd.south);
  \draw[->, thick, ln-accent] (rd.west) -- (p3.east)
    node[midway, below, font=\scriptsize, text=ln-accent] {\iflangja{更新}{update}};
\end{tikzpicture}

  \caption{Renaming \texttt{ADD R3, R1, R2}.  R1 and R2 are looked up in the
    RAT, the result is given the free physical register P8, and the RAT
    entry of R3, which held P3, is overwritten with P8.}
  \label{fig:l06-rat}
\end{figure}

The physical registers are exactly the versions of \cref{sec:l06-false}.
Every write creates a new version, and the RAT records which version is
current in program order, so a reader automatically uses the version
produced by its closest preceding writer.

\begin{example}
  \label{ex:l06-rename}
  We rename the running example.  Before I1, the RAT maps R1 to R7, the
  registers the example uses, to P1 to P7, and the other architectural
  registers to physical registers that the example does not use; P8, P9, \ldots\ are free and are
  allocated in increasing order.  \Cref{tab:l06-rename} lists each renamed instruction
  and the RAT entries after it.  I2 reads P8, the result of I1, because I1
  has already updated the entry of R3.  I3 reads P4, the old R4, before I4
  maps R4 to P11.  The two writes to R5 go to P9 and P12, and at the end the
  RAT maps R5 to P12, so an instruction that follows reads the result of I5,
  as program order requires.
\end{example}

\begin{table}[htbp]
  \centering\small
  \caption{Renaming the running example.  Entries changed by an instruction
    are in bold; R1, R2 and R7 keep P1, P2 and P7 throughout.}
  \label{tab:l06-rename}
  \begin{tabular}{@{}lllcccc@{}}
    \toprule
    & Instruction & Renamed & R3 & R4 & R5 & R6 \\
    \midrule
    & (initial) & & P3 & P4 & P5 & P6 \\
    I1 & \texttt{ADD R3, R1, R2} & \texttt{ADD P8, P1, P2}
       & \textbf{P8} & P4 & P5 & P6 \\
    I2 & \texttt{SUB R5, R3, R4} & \texttt{SUB P9, P8, P4}
       & P8 & P4 & \textbf{P9} & P6 \\
    I3 & \texttt{AND R6, R1, R4} & \texttt{AND P10, P1, P4}
       & P8 & P4 & P9 & \textbf{P10} \\
    I4 & \texttt{OR~~R4, R2, R7} & \texttt{OR~~P11, P2, P7}
       & P8 & \textbf{P11} & P9 & P10 \\
    I5 & \texttt{XOR R5, R1, R7} & \texttt{XOR P12, P1, P7}
       & P8 & P11 & \textbf{P12} & P10 \\
    \bottomrule
  \end{tabular}
\end{table}

\section{Dependences after renaming}
\label{sec:l06-after-renaming}

In a renamed program every result is written into its own physical
register.\source{Lesson 6, video 12; H\&P §3.4.}  No two instructions write
the same location, and no instruction writes a location that an earlier
instruction reads, so WAW and WAR dependences between register accesses no
longer exist.\margin{Dependences through memory locations are not affected
by register renaming; Lesson~9 treats them.}  Every read,
on the other hand, still refers to the result of the instruction that
produced the value in the original program, so every RAW dependence is
preserved.  In \cref{tab:l06-rename} the only dependence left is that of I2
on I1 through P8.

Same-cycle execution now only has to respect this one dependence.  I1, I3,
I4 and I5 execute in the first cycle and I2 in the second, with the result
of I1 forwarded to it.  I2 and I5 write P9 and P12, so the order of their
writes is irrelevant, and I2 reads P4, which I4 does not touch.  The five
instructions execute in two cycles, and every later reader still finds the
correct values through the RAT.

\begin{keypoint}
  RAW dependences are the only true dependences between instructions.
  Register renaming removes all WAW and WAR dependences on registers, so the
  only
  constraint left on executing instructions in parallel is the flow of data
  itself.
\end{keypoint}

\section{Instruction-level parallelism}
\label{sec:l06-ilp}

Once false dependences are removed, the number of cycles needed to execute a
program depends only on its RAW dependences.\source{Lesson 6, videos 13--15;
H\&P ch.~3; \cite{wall1991}.}  This number is a property of the program
itself, not of a particular processor.  It is measured on a hypothetical
\term{ideal processor}[理想プロセッサ], which
\begin{itemize}
  \item renames registers, so that no false dependences exist;
  \item executes every instruction in the first cycle in which all the
    values it reads are available: the first cycle if it depends on no
    earlier instruction, otherwise the cycle after the latest of the
    instructions it depends on;
  \item has no limit on the number of instructions it executes in a cycle;
  \item has unlimited hardware, so that no instruction waits for a busy
    execution unit: there are no structural dependences;
  \item executes every instruction in exactly one cycle, whatever its
    operation.
\end{itemize}

\begin{definition}[Instruction-level parallelism]
  \label{def:l06-ilp}
  The \term{instruction-level parallelism}[命令レベル並列性]%
  \index{ILP|see{instruction-level parallelism}} (ILP) of a program that
  executes $N$ instructions is
  \begin{equation}
    \text{ILP} \;=\; \frac{N}{C} ,
    \label{eq:l06-ilp}
  \end{equation}
  where $C$ is the number of cycles in which the ideal processor executes
  them.
\end{definition}

The cycle $c(I)$ in which the ideal processor executes instruction $I$
follows directly from the RAW dependences:
\begin{equation}
  c(I) \;=\; 1 + \max_{J \in D(I)} c(J) ,
  \qquad C = \max_{I} c(I) ,
  \label{eq:l06-cycle}
\end{equation}
where $D(I)$ is the set of instructions whose results $I$ reads, and the
maximum over an empty set is~0.  $C$ is thus the length of the longest chain
of RAW dependences in the program.  Since no processor can execute an
instruction before the values it needs exist, the ILP is an upper bound on the
rate, in instructions per cycle, at which any processor can execute this
program (\cref{sec:l06-ipc}).\tip{Renaming need not be done on paper: for
each register read, find the closest \emph{preceding} writer; later writes
do not matter.}

\begin{example}
  \label{ex:l06-ilp}
  We compute the ILP of the following eight instructions:
  \begin{lstlisting}[language={[x86masm]Assembler}, numbers=none]
I1:  ADD R1, R2, R3
I2:  MUL R4, R1, R5    ; R1 from I1; R5 written later
I3:  SUB R6, R2, R7    ; R2, R7 written only later
I4:  ADD R1, R6, R8    ; R6 from I3
I5:  AND R9, R1, R8    ; R1 from I4, not from I1
I6:  OR  R2, R8, R3
I7:  XOR R7, R4, R2    ; R4 from I2, R2 from I6
I8:  ADD R5, R3, R8
\end{lstlisting}
  I1, I3, I6 and I8 read only registers not written before them and execute
  in cycle~1.  I2 and I4 depend on instructions of cycle~1 and execute in
  cycle~2.  I5 and I7 each depend on an instruction of cycle~2 and execute
  in cycle~3 (\cref{fig:l06-ilp-graph}).  Hence $C=3$ and
  $\text{ILP} = 8/3 \approx 2.67$.  Two pitfalls are visible.  R1 is written
  by both I1 and I4, and I5 reads the value of I4; taking I1 as its producer
  would wrongly place I5 in cycle~2.  And I3 reads R2 and R7, which I6 and I7
  write only later; these are WAR dependences, removed by renaming, and they
  do not delay I3.
\end{example}

\begin{figure}[htbp]
  \centering
  % RAW dependences of the ILP example, instructions placed in the cycle in
% which the ideal processor executes them.
\begin{tikzpicture}[x=1cm,y=1cm, font=\footnotesize\sffamily, >={Stealth[length=1.8mm]},
  ins/.style={draw, fill=ln-box, rounded corners=2pt, align=center,
    minimum width=2.4cm, minimum height=0.75cm, inner sep=1pt,
    font=\scriptsize\sffamily},
  dep/.style={->, thick, ln-accent},
  reg/.style={font=\scriptsize\sffamily, fill=white, inner sep=1pt}]
  \foreach \c/\x in {1/0,2/3.9,3/7.8} {
    \node[font=\footnotesize\bfseries\sffamily, text=ln-muted] at (\x,0.75)
      {\iflangja{サイクル \c}{cycle \c}};
  }
  \draw[ln-rule] (1.95,0.95) -- (1.95,-3.4);
  \draw[ln-rule] (5.85,0.95) -- (5.85,-3.4);
  % cycle 1
  \node[ins] (i1) at (0,0)    {I1\\\texttt{ADD R1,R2,R3}};
  \node[ins] (i6) at (0,-1)   {I6\\\texttt{OR~~R2,R8,R3}};
  \node[ins] (i3) at (0,-2)   {I3\\\texttt{SUB R6,R2,R7}};
  \node[ins] (i8) at (0,-3)   {I8\\\texttt{ADD R5,R3,R8}};
  % cycle 2
  \node[ins] (i2) at (3.9,0)  {I2\\\texttt{MUL R4,R1,R5}};
  \node[ins] (i4) at (3.9,-2) {I4\\\texttt{ADD R1,R6,R8}};
  % cycle 3
  \node[ins] (i7) at (7.8,-0.5) {I7\\\texttt{XOR R7,R4,R2}};
  \node[ins] (i5) at (7.8,-2)   {I5\\\texttt{AND R9,R1,R8}};
  % RAW dependences, labelled with the register
  \draw[dep] (i1) -- node[reg] {R1} (i2);
  \draw[dep] (i3) -- node[reg] {R6} (i4);
  \draw[dep] (i4) -- node[reg] {R1} (i5);
  \draw[dep] (i2.east) -- node[reg] {R4} (i7.north west);
  \draw[dep] (i6.east) -- node[reg, pos=0.3] {R2} (i7.south west);
\end{tikzpicture}

  \caption{RAW dependences of \cref{ex:l06-ilp}, labelled with the register
    that carries the value.  Each instruction is placed in the cycle in which
    the ideal processor executes it; the longest chain has length~3.}
  \label{fig:l06-ilp-graph}
\end{figure}

\section{Structural and control dependences}
\label{sec:l06-control}

Structural dependences, in which an instruction waits for hardware that
another instruction is using, do not arise on the ideal processor, which
has enough hardware for every instruction that is ready.\source{Lesson 6,
video 16; H\&P §3.1.}  Branches, however, add control dependences: whether
the instructions after a branch are executed depends on its outcome.  The
ideal processor that defines the ILP is also assumed to predict every
branch perfectly.  It therefore knows from the start which instructions are
executed, and the instructions after a branch do not wait for the branch to
execute.  The branch itself is an instruction like any other: it counts
towards~$N$ and executes one cycle after the instructions that produce its
operands.  Instructions skipped because of the branch outcome are not
executed and do not count.

\begin{example}
  \label{ex:l06-control}
  In the following fragment the branch is taken, so I4 is not executed:
  \begin{lstlisting}[language={[x86masm]Assembler}, numbers=none]
I1:        ADD R1, R2, R3     ; cycle 1
I2:        SUB R4, R1, R5     ; cycle 2: R1 from I1
I3:        BEQ R4, R6, Skip   ; cycle 3: R4 from I2
I4:        ADD R7, R8, R9     ; not executed
I5:  Skip: XOR R8, R2, R6     ; cycle 1
I6:        AND R9, R8, R1     ; cycle 2: R8 from I5
\end{lstlisting}
  With perfect branch prediction I5 executes in cycle~1 although the branch
  executes only in cycle~3.  Five instructions execute in 3 cycles, so
  $\text{ILP} = 5/3 \approx 1.67$.  If I5 and I6 had to wait for the branch,
  they would execute in cycles 4 and~5, and the ILP would drop to~1.
\end{example}

\section{ILP and IPC}
\label{sec:l06-ipc}

The ILP describes a program on the ideal processor; a real processor is
described by the rate it actually achieves.\source{Lesson 6, videos 17--20;
H\&P ch.~3.}

\begin{definition}[Instructions per cycle]
  \label{def:l06-ipc}
  The average number of instructions a processor completes per cycle is its
  \term{instructions per cycle}[サイクル当たり命令数]%
  \index{IPC|see{instructions per cycle}} (IPC).  It is the reciprocal of
  the CPI of the iron law (\cref{eq:iron-law}): $\text{IPC} =
  1/\text{CPI}$.
\end{definition}

For every processor and program, $\text{IPC} \le \text{ILP}$, because the
ideal processor executes each instruction as early as the dependences allow.
A real processor falls short of the ideal for several reasons.
\begin{description}
  \item[Width] It can execute only a limited number $w$ of instructions per
    cycle.  A \emph{narrow} processor executes one or two, a \emph{wide} one
    four or more.  Whatever the program, $\text{IPC} \le w$.
  \item[Execution units] It has a limited number of adders, multipliers and
    other execution units, and operations such as multiplication may take
    several cycles.  An instruction whose kind of unit is busy waits even
    when its inputs are ready: a structural dependence, which the ideal
    processor does not have.
  \item[Execution order] A processor with \term{in-order
    execution}[インオーダ実行] never executes an instruction before an
    instruction that precedes it in the program.  Consecutive independent
    instructions may execute in the same cycle, but an instruction that has
    to wait for its inputs also holds back every instruction behind it.  A
    processor with \emph{out-of-order}\index{out-of-order execution}
    execution may execute any instruction whose inputs are available
    (Lesson~7).
\end{description}
An out-of-order processor needs renaming to do so; without it, it would
also have to respect the WAW and WAR dependences of
\cref{sec:l06-dependences}.  An in-order processor reads and writes
registers in program order and needs no renaming.

\begin{example}
  \label{ex:l06-ipc}
  We run the program of \cref{ex:l06-ilp} on four processors, each
  executing every instruction in one cycle and having enough execution units
  of every kind; the out-of-order ones rename
  registers and, when more instructions are ready than the width allows,
  execute the oldest ones first.  \Cref{tab:l06-ipc} lists the instructions
  executed in each cycle.  The in-order processors execute only I1 in
  cycle~1, because I2 must wait for R1 and blocks I3 to I8.  The same
  blocking recurs at I4 (waiting for R6), at I5 (waiting for R1 from I4) and
  at I7 (waiting for R2), so these processors need 5 cycles whatever their
  width.  The two-wide out-of-order processor executes two
  ready instructions in every cycle and needs 4 cycles; only the four-wide
  out-of-order processor reaches the ILP.
\end{example}

\begin{table}[htbp]
  \centering\small
  \caption{Instructions of \cref{ex:l06-ilp} executed in each cycle.  The
    schedule of the ideal processor equals that of the four-wide
    out-of-order processor.}
  \label{tab:l06-ipc}
  \begin{tabular}{@{}lcccccc@{}}
    \toprule
    Processor & Cycle 1 & Cycle 2 & Cycle 3 & Cycle 4 & Cycle 5 & IPC \\
    \midrule
    In-order, 2 wide     & 1          & 2, 3 & 4    & 5, 6 & 7, 8 & 1.60 \\
    In-order, 4 wide     & 1          & 2, 3 & 4    & 5, 6 & 7, 8 & 1.60 \\
    Out-of-order, 2 wide & 1, 3       & 2, 4 & 5, 6 & 7, 8 &      & 2.00 \\
    Out-of-order, 4 wide & 1, 3, 6, 8 & 2, 4 & 5, 7 &      &      & 2.67 \\
    \bottomrule
  \end{tabular}
\end{table}

In the example both in-order processors are held back by the order:
widening from two to four gains nothing.  Out-of-order execution lifts the
two-wide processor only to its width (from 1.60 to $w = 2$), but the
four-wide one to the ILP (from 1.60 to 2.67).  In general
(\cref{tab:l06-limits}), the width is the main limit of a narrow processor
(with $w = 1$, $\text{IPC} \le 1$ in any order), so out-of-order execution
raises its IPC only moderately.  In a wide in-order processor, waiting instructions
block the ones behind them and most of the width stays unused.  Wide
processors therefore need out-of-order execution to use their width, and
out-of-order execution pays off fully only in a wide processor.

\begin{table}[htbp]
  \centering\small
  \caption{What limits the IPC of a processor relative to the ILP of the
    program.}
  \label{tab:l06-limits}
  \begin{tabular}{@{}lll@{}}
    \toprule
    & In-order & Out-of-order \\
    \midrule
    Narrow (1--2 wide) & mainly width & mainly width \\
    Wide (4 or more)   & mainly order; width unused & little: IPC approaches ILP \\
    \bottomrule
  \end{tabular}
\end{table}

\begin{keypoint}[Lesson summary]
  Executing many instructions in the same cycle would bring the CPI far below
  1, but data dependences prevent it.  Only RAW dependences are true
  dependences; WAW and WAR dependences come from reusing register names.
  Register renaming, driven by the register alias table, maps architectural
  registers to newly allocated physical registers and removes all false
  dependences between register accesses.  The ILP (\cref{eq:l06-ilp}) of a
  program is the number of instructions per cycle of an ideal processor with
  renaming, unlimited width, no structural dependences and perfect branch
  prediction: the instruction count divided by the
  length of the longest RAW chain; it is a property of the program and bounds
  the IPC of every processor.  Real processors fall below it, limited by
  their width, their execution units and in-order execution; approaching the ILP takes a wide
  out-of-order processor.  The next lesson presents Tomasulo's
  algorithm, which executes instructions out of order with renaming.
\end{keypoint}

\end{document}
